\chapter{The Jagged Frontier}
\label{ch:jagged-frontier}

\epigraph{The person who knows where the frontier is jagged in their domain is the most valuable person in the organisation.}{}

\section{Why Wholesale Automation Is Reckless}

If coordination is now computationally cheap, the logical next question is whether to simply replace the hierarchy with AI. Several prominent voices (most notably \textcite{DorseyBotha2026}) propose exactly this. But the evidence reveals that the boundary between what AI can do and what it cannot is dangerously unpredictable. This boundary is the single most important concept in the entire coordination collapse thesis.

\section{The BCG Experiment}

The \meshterm{jagged technological frontier} was identified and named by \textcite{DellAcqua2026} in what remains the most rigorous field experiment on human-AI collaboration in knowledge work. The study was preregistered, conducted with 758 consultants at Boston Consulting Group: not a laboratory convenience sample, not a blog post, but a controlled experiment with one of the world's premier professional services firms.

The researchers assigned consultants to eighteen realistic consulting tasks and measured the results against a control group working without AI assistance. The findings established two sharply distinct performance regimes.

\begin{evidencebox}[title={Inside the Frontier}]
For the eighteen tasks that fell within AI capability:
\begin{itemize}
  \item Consultants using GPT-4 completed \textbf{12.2 per cent more tasks}
  \item They finished \textbf{25.1 per cent faster}
  \item They produced \textbf{40 per cent higher quality} output, as rated by blind evaluators
\end{itemize}
\end{evidencebox}

\begin{cautionbox}[title={Outside the Frontier}]
For a complex business judgement task requiring integration of real-world context:
\begin{itemize}
  \item AI users were \textbf{19 per cent less likely} to produce correct solutions than those working without AI
  \item The frontier was unpredictable: adjacent tasks that appeared similar fell on opposite sides
\end{itemize}
\end{cautionbox}

The metaphor of a ``jagged'' frontier is precise. Unlike a smooth gradient where tasks become progressively harder for AI, the boundary is irregular and counterintuitive. Two tasks that appear nearly identical from the outside (similar in domain, complexity, and required expertise) can fall on opposite sides of the frontier. One is dramatically improved by AI assistance; the other is degraded by it. There is no reliable way to determine which is which without empirical testing.

\begin{figure}[h]
\centering
\includegraphics[width=0.92\textwidth]{nb-jagged-frontier.jpg}
\caption{The jagged technological frontier. Adjacent tasks fall on opposite sides of an irregular, moving boundary; the only reliable map is empirical testing in your own domain.}
\label{fig:jagged-frontier}
\end{figure}

\section{The Skill-Levelling Effect}

The BCG experiment revealed a second finding with profound implications for organisational design. AI did not improve all consultants equally:

\begin{itemize}
  \item Bottom-half performers saw a \textbf{43 per cent performance jump} with AI assistance.
  \item Top performers improved measurably but by a smaller margin.
  \item AI compressed the talent distribution, bringing lower-performing consultants much closer to the performance of their higher-performing colleagues.
\end{itemize}

This skill-levelling effect changes what managers are managing. In a pre-AI organisation, individual performance variation is a primary concern: identifying, developing, and retaining high performers is a core management function. When AI compresses the performance distribution, the variation between individuals shrinks. The manager's value shifts from talent differentiation to system coherence, ensuring that the newly capable workforce is pointed at the right problems and that AI assistance is applied where it helps rather than where it harms.

\section{Three Collaboration Modes}

A follow-up study by \textcite{Randazzo2024}, conducted with 244 BCG consultants, identified three distinct modes of human-AI collaboration. These modes are not merely descriptive; they are predictive of performance outcomes.

\begin{table}[h]
\centering
\caption{Human-AI collaboration modes and performance outcomes}
\label{tab:collaboration-modes}
\begin{tabularx}{\textwidth}{l L{5cm} X}
\toprule
\textbf{Mode} & \textbf{Description} & \textbf{Performance} \\
\midrule
Centaurs & Clear division of labour. The human performs tasks where human judgment is strongest; the AI handles tasks inside the frontier. Distinct handoff points between human and machine. & High \\
\addlinespace
Cyborgs & Continuous interleaving. Human and AI work on the same analysis, passing work back and forth across the frontier boundary. No clean separation of responsibilities. & High \\
\addlinespace
Self-Automators & Wholesale delegation to AI with minimal oversight. The human provides the initial prompt and accepts the output with little or no review. & Lowest \\
\bottomrule
\end{tabularx}
\end{table}

Both Centaurs and Cyborgs outperformed solo humans and solo AI. Self-Automators produced the weakest results, precisely because wholesale delegation triggers the vigilance problem examined in the next chapter.

\begin{keyinsight}[title={The Organisational Design Implication}]
This taxonomy maps directly onto the three organisational models examined in Chapter~\ref{ch:three-models}:
\begin{itemize}
  \item Block's intelligence layer model creates \textbf{Self-Automator organisations}, delegating coordination to an AI world model.
  \item Every's agent colony model creates \textbf{Cyborg organisations}: continuous interleaving between human and personal agent.
  \item The Dynamic Agentic Mesh creates \textbf{Centaur organisations}: clear boundaries with judgment brokers at the handoff points.
\end{itemize}
The evidence shows that Self-Automators perform worst. This is the empirical case against wholesale delegation.
\end{keyinsight}

\section{Quantifying the Frontier: The Remote Labor Index}

The BCG experiment and its successors establish that the frontier is jagged. They do not, on their own, tell us how fast the frontier is moving, or how far it has moved. That longitudinal, quantitative measurement now exists, and it changes the shape of the argument.

The \meshterm{Remote Labor Index} (RLI) \parencite{RLI2025}, published by the Center for AI Safety and Scale AI as arXiv:2510.26787 in October 2025, tests AI models against real, economically valuable freelance projects: 3D modelling and CAD, architecture, graphic design, video and audio editing, data analysis, and web-application programming. Every project is judged by human evaluators against a gold-standard deliverable produced by a paid professional. This is a materially harsher bar than the task-level benchmarks that dominate AI capability reporting. The RLI does not ask whether a model can pass a coding exercise or answer a reasoning puzzle; it asks whether the model's output would be accepted by a paying client in place of a professional's work.

At launch, the best-performing model automated 2.5 per cent of RLI projects to that standard. The July 2026 update \parencite{CAIS2026RLI} reports a step change: Claude Fable 5 at 16.1 per cent, Claude Opus 4.8 at 8.3 per cent, and GPT-5.5 at 6.3 per cent. In the Center for AI Safety's own words: ``The frontier has more than quadrupled in under eight months, a concrete signal of how quickly economically capable AI agents are advancing.''

The update's example deliverables illustrate what 16.1 per cent automation looks like in practice: photorealistic renders of an engagement ring redesign, a flat-design 2D animated advertisement for a fictional tree-services company, and dimensioned floor plans with accompanying photorealistic renders. These are not toy tasks. They are the kind of whole-deliverable freelance work that supports real careers. CAIS is explicit about the limits of what has been achieved: ``Today's AI still falls short of professional quality on most projects.''

\begin{keyinsight}[title={The First Longitudinal Quantification of the Jagged Frontier}]
The RLI is the first instrument that tracks the jagged frontier longitudinally, at client-quality standard, across whole deliverables rather than sub-tasks. It deserves to be read two ways at once, and the report insists on holding both readings simultaneously rather than collapsing into either.

\textbf{Read one:} 16.1 per cent means the human remains indispensable on roughly five of every six projects. The gold-standard bar (would a paying client accept this?) is not being cleared for the overwhelming majority of real freelance work. Professional-quality completion is still the exception, not the rule.

\textbf{Read two:} the boundary moved more than sixfold in eight months. A frontier that automates one project in forty in October 2025 and one project in six by July 2026 is not moving on the multi-year timescale that organisational change management assumes. It is moving on a timescale shorter than most transformation programmes.

Both readings are correct, and the tension between them is not a contradiction to be resolved; it is the terrain the judgment broker (Chapter~\ref{ch:role-transformation}) operates on. Tasks are not jobs: automating 16.1 per cent of freelance \emph{projects} at professional quality is a different claim from automating 16.1 per cent of freelance \emph{work}, let alone 16.1 per cent of jobs. The RLI measures whole deliverables against a client-acceptance bar, which makes it a better proxy for job-level displacement than task-level benchmarks, but it is still a proxy, not a labour-market outcome. The space between ``the human is still needed five times out of six'' and ``the boundary moved sixfold in eight months'' is exactly the space that frontier-mapping expertise exists to navigate. Knowing where today's one-in-six lives, and how fast it is becoming one-in-five, is the scarce and valuable knowledge this report has been calling the judgment broker's domain since the opening chapters.
\end{keyinsight}

The RLI is a measure of capability under harsh, client-quality evaluation. A complementary measure (of capability under a looser, time-based standard) comes from METR's task-horizon work. METR's metric asks how long a task a frontier model can complete with 50 per cent reliability, calibrated against the time a human expert would need. That horizon has doubled roughly every seven months since 2019, and the trend has recently accelerated \parencite{METR2026}. Claude Opus 4.6 was measured at approximately 14.5 human-expert hours when added to METR's tracker in February 2026, meaning frontier models now reliably complete, half the time, tasks that would occupy a skilled professional for nearly two working days.

Both metrics point the same direction and both come with the same caveat. These are benchmark capabilities, measured under conditions designed to isolate the model's raw ability. Deployed capability is narrower: real-world use is gated by system integration, permission scoping, and the handling of exceptions that a benchmark does not surface. The RLI and METR describe how fast the frontier is moving in principle. Chapter~\ref{ch:vigilance} and the mesh architecture in later chapters describe what it takes to move it safely in practice.

\section{Why the Frontier Matters for Governance}

The jagged frontier has four implications that recur throughout this report.

\textbf{First, it invalidates wholesale automation.} If you cannot predict in advance which tasks AI will improve and which it will degrade, a strategy of replacing human judgment with AI judgment is not bold --- it is uninformed. The 19 per cent performance degradation on tasks outside the frontier is not a temporary limitation that will be resolved by the next model generation. It is a structural feature of deploying any system at the boundary of its capability.

\textbf{Second, it defeats one-time change playbooks.} The frontier moves as models improve. A task that falls outside the frontier with GPT-4 may fall inside it with the next generation. Conversely, increased model capability can create new failure modes at different points on the boundary. A change management plan that maps the frontier today will be outdated within months. This demands continuous adaptation, not phased transformation.

\textbf{Third, it creates a new form of expertise.} The person who knows where the frontier falls in their specific domain (who has mapped which tasks are reliably improved by AI and which are reliably degraded) holds knowledge that is genuinely scarce and strategically valuable. This is the \meshterm{judgment broker} role developed in Chapter~\ref{ch:role-transformation}. It is not a euphemism for a diminished middle management function. It is a description of expertise that did not exist before the frontier existed.

\textbf{Fourth, it demands active engagement, not passive oversight.} The standard response to AI risk (``keep a human in the loop'') assumes that the human will maintain vigilance. The next chapter demonstrates why that assumption fails, and what it means for organisational design.

\section{The Frontier Is Dynamic}

It is important to note a limitation of the evidence. The BCG experiment used GPT-4, now several model generations behind current frontier models. The specific task boundaries identified in 2023 have almost certainly shifted. Some tasks that fell outside the frontier are now inside it; others may have developed new failure modes as model capabilities have changed. The April 2026 edition of this report flagged this as an open problem: the BCG task map was already outdated, and no instrument existed to re-measure the frontier as models improved.

That instrument now exists. The Remote Labor Index, tracked from its October 2025 launch through the July 2026 update, is exactly the longitudinal re-measurement this chapter previously had to do without. It confirms both halves of the original claim at once. The frontier moves fast: a sixfold improvement in eight months is not a rounding error, it is a different pace of change than most organisational planning horizons assume. And the frontier remains jagged: at 16.1 per cent automation, professional-quality completion of a whole freelance deliverable is still the exception rather than the rule, across every category the index measures. The GPT-4-era task map is out of date, as flagged; what has changed since is that the report no longer has to rely on a single snapshot to make the point. Continuous measurement (the RLI updated every few months, METR's tracker updated as models are released) is now doing the job that a one-off field experiment could only gesture at.

The implication is unchanged, and sharper for it: organisations need people who actively probe the boundary rather than assuming they know where it is. The frontier will always be jagged because AI capability will always be unevenly distributed across the task space of any complex domain. What is new is that the rate of movement is now measured, not merely inferred, and the rate is fast enough that a frontier map drawn this quarter should be assumed stale by the next.

The publication of the BCG study in \emph{Organization Science} in March 2026 \parencite{DellAcqua2026} provides the peer-reviewed validation that the working paper lacked. The jagged frontier is no longer a provocative hypothesis. It is an established finding in the organisational science literature --- and, as of the Remote Labor Index, a finding with a running clock attached to it.
